% GENERATED FILE. Edit canonical lessons, then run npm run generate:books.
% canonical-source-hash: fnv1a64:15f3804083e61a34
% canonical-lessons: BN-C341-longghon-kora, BN-C341-omanyo-kora, BN-C341-jaliyati, BN-C341-durniti, BN-C341-onyay

\chapter{Bending the Rules}
\label{ch:bn-bending-the-rules}
\hlchaptermodality{341}

\begin{chapteropening}
\textbf{By the end of this chapter:} \emph{I can say when a rule was broken and what was unfair.}

Complete the last lesson of chapter 341: I can say when a rule was broken and what was unfair.
\end{chapteropening}

\section[\texorpdfstring{\bn{লঙ্ঘন} \bn{করা} (lôngghôn kôrā)}{লঙ্ঘন করা (lôngghôn kôrā)}]{\bn{লঙ্ঘন} \bn{করা} (lôngghôn kôrā) — to violate; to break (a law, an agreement)}

\label{lesson:BN-C341-longghon-kora}



\begin{quote}
Before the new one: say the Bengali for to carry out, then the Bengali for to enforce.
\end{quote}

\subsection*{You'll want to know: \bn{লঙ্ঘন} \bn{করা}}
\textbf{\bn{লঙ্ঘন} \bn{করা}} — \emph{lôngghôn kôrā} — \textquotedblleft{}to violate; to break (a law, an agreement)\textquotedblright{}.

\begin{quote}
\bn{কোম্পানিটা} \bn{চুক্তি} \bn{লঙ্ঘন} \bn{করল।}
\end{quote}

\begin{quote}
\emph{kompāniṭā chukti lôngghôn kôrlô.}
\end{quote}

\begin{quote}
The company broke the contract.
\end{quote}

Past (he/she …): \bn{তিনি} \bn{নিয়মটা} \bn{লঙ্ঘন} \bn{করলেন}, \emph{tini niyômṭā lôngghôn kôrlen}.

\begin{cousinweb}
\textbf{\bn{লঙ্ঘন} \bn{করা}} (\emph{lôngghôn kôrā}) is Sanskrit \emph{laṅghana}, leaping over and so overstepping, a violation (Monier-Williams), with \bn{করা}: to step across the line a law draws.
\end{cousinweb}

\begin{grammarlens}[title={What you've built}]
One more word for giving reasons and opinions.
\end{grammarlens}

\subsection*{Your turn}
\begin{itemize}
\raggedright
  \item \emph{Say it:} \emph{lôngghôn kôrā}
  \item \emph{Say it:} \emph{lôngghôn kôrā}, once more
  \item \emph{Say it:} \emph{kompāniṭā chukti lôngghôn kôrlô.} — The company broke the contract.
  \item \emph{Recall:} say the Bengali for to carry out, then the Bengali for to enforce, then say \emph{lôngghôn kôrā} again
\end{itemize}

\begin{tcolorbox}[breakable,colback=teal!4,colframe=teal!35!black,title={Before you move on}]
What does \bn{লঙ্ঘন} \bn{করা} mean? (\textquotedblleft{}To violate; to break (a law, an agreement)\textquotedblright{}.) Say it once more.
\end{tcolorbox}
\section[\texorpdfstring{\bn{অমান্য} \bn{করা} (ômānyô kôrā)}{অমান্য করা (ômānyô kôrā)}]{\bn{অমান্য} \bn{করা} (ômānyô kôrā) — to disobey; to defy}

\label{lesson:BN-C341-omanyo-kora}



\begin{quote}
Before the new one: say the Bengali for to enforce, then the Bengali for to violate.
\end{quote}

\subsection*{You'll want to know: \bn{অমান্য} \bn{করা}}
\textbf{\bn{অমান্য} \bn{করা}} — \emph{ômānyô kôrā} — \textquotedblleft{}to disobey; to defy\textquotedblright{}.

\begin{quote}
\bn{ছাত্ররা} \bn{নিয়ম} \bn{অমান্য} \bn{করল।}
\end{quote}

\begin{quote}
\emph{chhātrôrā niyôm ômānyô kôrlô.}
\end{quote}

\begin{quote}
The students disobeyed the rules.
\end{quote}

Past (he/she …): \bn{তিনি} \bn{আইন} \bn{অমান্য} \bn{করলেন}, \emph{tini āin ômānyô kôrlen}.

\begin{cousinweb}
\textbf{\bn{অমান্য} \bn{করা}} (\emph{ômānyô kôrā}) is Sanskrit \emph{a-} 'not' with \emph{mānya}, 'to be respected' (Monier-Williams), and \bn{করা}: to treat a rule as not worth heeding, on purpose.
\end{cousinweb}

\begin{grammarlens}[title={What you've built}]
One more word for giving reasons and opinions.
\end{grammarlens}

\subsection*{Your turn}
\begin{itemize}
\raggedright
  \item \emph{Say it:} \emph{ômānyô kôrā}
  \item \emph{Say it:} \emph{ômānyô kôrā}, once more
  \item \emph{Say it:} \emph{chhātrôrā niyôm ômānyô kôrlô.} — The students disobeyed the rules.
  \item \emph{Recall:} say the Bengali for to enforce, then the Bengali for to violate, then say \emph{ômānyô kôrā} again
\end{itemize}

\begin{tcolorbox}[breakable,colback=teal!4,colframe=teal!35!black,title={Before you move on}]
What does \bn{অমান্য} \bn{করা} mean? (\textquotedblleft{}To disobey; to defy\textquotedblright{}.) Say it once more.
\end{tcolorbox}
\section[\texorpdfstring{\bn{জালিয়াতি} (jāliyāti)}{জালিয়াতি (jāliyāti)}]{\bn{জালিয়াতি} (jāliyāti) — forgery; fraud}

\label{lesson:BN-C341-jaliyati}



\begin{quote}
Before the new one: say the Bengali for to violate, then the Bengali for to disobey.
\end{quote}

\subsection*{You'll want to know: \bn{জালিয়াতি}}
\textbf{\bn{জালিয়াতি}} — \emph{jāliyāti} — \textquotedblleft{}forgery; fraud\textquotedblright{}.

\begin{quote}
\bn{পুলিশ} \bn{একটা} \bn{বড়} \bn{জালিয়াতি} \bn{ধরল।}
\end{quote}

\begin{quote}
\emph{pulish ekṭā bôṛo jāliyāti dhôrlô.}
\end{quote}

\begin{quote}
The police uncovered a big fraud.
\end{quote}

\begin{cousinweb}
\textbf{\bn{জালিয়াতি}} (\emph{jāliyāti}) is built on \bn{জাল}, 'counterfeit', usually traced to Arabic \emph{jaʿl}, fabrication, through Persian; it is a different word from \bn{জাল} 'a net'. The -\bn{ইয়াতি} ending makes the act.
\end{cousinweb}

\begin{grammarlens}[title={What you've built}]
One more word for giving reasons and opinions.
\end{grammarlens}

\subsection*{Your turn}
\begin{itemize}
\raggedright
  \item \emph{Say it:} \emph{jāliyāti}
  \item \emph{Say it:} \emph{jāliyāti}, once more
  \item \emph{Say it:} \emph{pulish ekṭā bôṛo jāliyāti dhôrlô.} — The police uncovered a big fraud.
  \item \emph{Recall:} say the Bengali for to violate, then the Bengali for to disobey, then say \emph{jāliyāti} again
\end{itemize}

\begin{tcolorbox}[breakable,colback=teal!4,colframe=teal!35!black,title={Before you move on}]
What does \bn{জালিয়াতি} mean? (\textquotedblleft{}Forgery; fraud\textquotedblright{}.) Say it once more.
\end{tcolorbox}
\section[\texorpdfstring{\bn{দুর্নীতি} (durnīti)}{দুর্নীতি (durnīti)}]{\bn{দুর্নীতি} (durnīti) — corruption}

\label{lesson:BN-C341-durniti}



\begin{quote}
Before the new one: say the Bengali for to disobey, then the Bengali for forgery.
\end{quote}

\subsection*{You'll want to know: \bn{দুর্নীতি}}
\textbf{\bn{দুর্নীতি}} — \emph{durnīti} — \textquotedblleft{}corruption\textquotedblright{}.

\begin{quote}
\bn{দুর্নীতির} \bn{বিরুদ্ধে} \bn{সবাই} \bn{একমত।}
\end{quote}

\begin{quote}
\emph{durnītir biruddhe sôbāi ekmôt.}
\end{quote}

\begin{quote}
Everyone is united against corruption.
\end{quote}

\begin{cousinweb}
\textbf{\bn{দুর্নীতি}} (\emph{durnīti}) is Sanskrit \emph{durnīti}, \emph{dur-} 'bad' with \emph{nīti} 'conduct, policy': bad government (Monier-Williams). Bengali narrowed it to bribes and abuse of office.
\end{cousinweb}

\begin{grammarlens}[title={What you've built}]
One more word for giving reasons and opinions.
\end{grammarlens}

\subsection*{Your turn}
\begin{itemize}
\raggedright
  \item \emph{Say it:} \emph{durnīti}
  \item \emph{Say it:} \emph{durnīti}, once more
  \item \emph{Say it:} \emph{durnītir biruddhe sôbāi ekmôt.} — Everyone is united against corruption.
  \item \emph{Recall:} say the Bengali for to disobey, then the Bengali for forgery, then say \emph{durnīti} again
\end{itemize}

\begin{tcolorbox}[breakable,colback=teal!4,colframe=teal!35!black,title={Before you move on}]
What does \bn{দুর্নীতি} mean? (\textquotedblleft{}Corruption\textquotedblright{}.) Say it once more.
\end{tcolorbox}
\section[\texorpdfstring{\bn{অন্যায়} (ônyāy)}{অন্যায় (ônyāy)}]{\bn{অন্যায়} (ônyāy) — unfair; an injustice}

\label{lesson:BN-C341-onyay}



\begin{quote}
Before the new one: say the Bengali for forgery, then the Bengali for corruption.
\end{quote}

\subsection*{You'll want to know: \bn{অন্যায়}}
\textbf{\bn{অন্যায়}} — \emph{ônyāy} — \textquotedblleft{}unfair; an injustice\textquotedblright{}.

\begin{quote}
\bn{এই} \bn{শাস্তিটা} \bn{অন্যায়।}
\end{quote}

\begin{quote}
\emph{ei shāstiṭā ônyāy.}
\end{quote}

\begin{quote}
This punishment is unfair.
\end{quote}

\begin{cousinweb}
\textbf{\bn{অন্যায়}} (\emph{ônyāy}) is Sanskrit \emph{anyāya}, \emph{a-} 'not' with \emph{nyāya} 'what is right': unjust action (Monier-Williams). It works as a noun, a wrong, and as an adjective, unfair.
\end{cousinweb}

\begin{grammarlens}[title={What you've built}]
One more word for giving reasons and opinions.
\end{grammarlens}

\subsection*{Your turn}
\begin{itemize}
\raggedright
  \item \emph{Say it:} \emph{ônyāy}
  \item \emph{Say it:} \emph{ônyāy}, once more
  \item \emph{Say it:} \emph{ei shāstiṭā ônyāy.} — This punishment is unfair.
  \item \emph{Recall:} say the Bengali for forgery, then the Bengali for corruption, then say \emph{ônyāy} again
\end{itemize}

\begin{tcolorbox}[breakable,colback=teal!4,colframe=teal!35!black,title={Before you move on}]
What does \bn{অন্যায়} mean? (\textquotedblleft{}Unfair; an injustice\textquotedblright{}.) Say it once more.
\end{tcolorbox}
